% ============================================================================
\chapter{管道与重定向}
\label{ch:pipe}
% Stage F-3 ~ F-4
% ============================================================================

\section{本章目标}

\begin{itemize}
  \item 理解 Unix 管道原理与进程间通信（IPC）分类
  \item 设计可插拔 \texttt{ipc\_ops\_t} 接口（匿名管道/命名管道/共享内存可扩展）
  \item 实现匿名管道后端（内核环形缓冲区 + 阻塞读写）
  \item 实现 \texttt{dup2()} 文件描述符复制
  \item Shell 集成管道 \texttt{|} 和重定向 \texttt{>}
\end{itemize}

% ============================================================================
\section{概念：进程间通信}
\label{sec:pipe-concept}
% ============================================================================

% TODO:
% - IPC 分类概述：
%     管道（无名/有名）、消息队列、共享内存、信号、Socket
% - 管道是最简单的 IPC：单向字节流，通过 fd 读写
% - 匿名管道（pipe）：父子进程间共享（fork 后继承 fd）
%     pipe(fds) → fds[0]=读端, fds[1]=写端
% - 命名管道（FIFO）：通过文件系统路径访问（mkfifo），不限于父子
% - 管道的 VFS 集成：管道也是 file 对象，读写通过 fd
% - 图：两个进程通过管道通信（进程A write → buffer → 进程B read）

\subsection{dup2 与 I/O 重定向}

% TODO:
% - dup2(oldfd, newfd) 语义：让 newfd 指向 oldfd 同一个 file 对象
% - Shell 重定向实现：
%     cmd > file  →  open(file) → dup2(fd, 1) → exec(cmd)
%     cmd < file  →  open(file) → dup2(fd, 0) → exec(cmd)
% - Shell 管道实现：
%     cmd1 | cmd2  →  pipe(fds) → fork → 子1: dup2(fds[1],1) → 子2: dup2(fds[0],0)

% ============================================================================
\section{可插拔接口：\texttt{ipc\_ops\_t}}
\label{sec:pipe-interface}
% ============================================================================

% TODO:
% - 接口设计：
%     typedef struct ipc_ops {
%         const char* name;                          // "pipe"
%         int (*create)(struct file* read_end, struct file* write_end);
%         int (*read)(struct file* f, void* buf, size_t count);
%         int (*write)(struct file* f, const void* buf, size_t count);
%         int (*close)(struct file* f);
%     } ipc_ops_t;
% - 为什么可插拔：未来可扩展 shared memory、message queue 等 IPC

% ============================================================================
\section{匿名管道后端}
\label{sec:pipe-impl}
% ============================================================================

% TODO:
% - 内核环形缓冲区：
%     struct pipe_buf { char data[PIPE_SIZE]; uint32_t read_pos, write_pos, count; }
%     PIPE_SIZE = 4096（一页）
% - write：如果 buffer 满 → 阻塞写端进程（BLOCKED）→ 等待读端读取
% - read：如果 buffer 空 → 阻塞读端进程 → 等待写端写入
% - EOF：所有写端关闭后 → read 返回 0
% - SIGPIPE：所有读端关闭后 → write 收到 SIGPIPE
% - dup2 实现：fd_table[newfd] = fd_table[oldfd]; refcount++

% ============================================================================
\section{验证}
\label{sec:pipe-verify}
% ============================================================================

% TODO:
% - 测试 1：pipe → fork → 父 write → 子 read → 匹配
% - 测试 2：dup2(fd, 1) → write(1, ...) → 输出到 fd 指向的文件
% - 测试 3：Shell 管道 echo "hello" | cat → 输出 hello
% - 测试 4：Shell 重定向 echo "data" > /tmp/test → 文件内容正确
% - shell 命令：pipe test

% ============================================================================
\section{本章小结}
% ============================================================================

管道通过可插拔 \texttt{ipc\_ops\_t} 接口实现进程间单向通信。
配合 \texttt{dup2} 的 fd 重定向能力，Shell 获得了管道 \texttt{|} 和重定向 \texttt{>} 的支持。
下一章实现 \texttt{mmap}，让用户态能直接映射内存和文件到地址空间。
